\documentclass[10pt]{article}
\usepackage[T1]{fontenc}
\usepackage[utf8]{inputenc}
\usepackage[french]{babel}
\usepackage{amsmath,amssymb,mathrsfs,mathpazo}
\usepackage{geometry}
\geometry{a4paper,margin=2cm}
\usepackage{graphicx,array,tabularx,tabularray,colortbl,multirow,multicol,booktabs,enumitem,pifont,ulem,fancybox,diagbox,cancel}
\usepackage{pstricks,pstricks-add,pst-plot,pst-tree,pst-node,pst-text,pst-func,pst-math,pst-bspline,pst-3dplot,pst-eucl,pst-solides3d}
\let\pscalculate\undefined
\usepackage{pst-bezier}
\usepackage{tikz,pgfplots,tkz-tab}
\usetikzlibrary{arrows,arrows.meta,patterns,calc,positioning,trees,decorations.markings,shapes.geometric,shapes.misc,shapes.arrows,intersections,decorations.pathmorphing,decorations.pathreplacing,backgrounds,shadows,babel}
\pgfplotsset{compat=1.18}
\usepgfplotslibrary{fillbetween}
\usepackage[np]{numprint}
\usepackage{listings,makecell,pas-tableur,twoopt}
\let\barre\undefined
\usepackage[pstricks]{tablvar}
\definecolor{darkgreen}{RGB}{0,100,0}
\providecommand{\red}{\color{red}}
\providecommand{\blue}{\color{blue}}
\providecommand{\green}{\color{green}}
\providecommand{\black}{\color{black}}
\providecommand{\white}{\color{white}}
\providecommand{\gray}{\color{gray}}
\providecommand{\cyan}{\color{cyan}}
\providecommand{\magenta}{\color{magenta}}
\providecommand{\yellow}{\color{yellow}}
\providecommand{\violet}{\color{violet}}
\providecommand{\orange}{\color{orange}}
\providecommand{\euro}{\texteuro}
\providecommand{\EUR}{\texteuro}
\providecommand{\ieme}{\textsuperscript{e}}
\providecommand{\degres}{\ensuremath{^\circ}}
\providecommand{\degree}{\ensuremath{^\circ}}
\providecommand{\vv}[1]{\overrightarrow{#1}}
\providecommand{\Coord}[1]{#1}
\providecommand{\decosix}{}
\providecommand{\starredbullet}{\ensuremath{\bullet}}
\pagestyle{empty}
\setlength{\parindent}{0pt}
\usepackage{adjustbox}\DeclareUnicodeCharacter{2236}{:}\DeclareUnicodeCharacter{2190}{\ensuremath{\leftarrow}}
\providecommand{\R}{\mathbb{R}}
\providecommand{\N}{\mathbb{N}}
\providecommand{\Z}{\mathbb{Z}}
\providecommand{\Q}{\mathbb{Q}}
\providecommand{\D}{\mathbb{D}}
\providecommand{\C}{\mathbb{C}}
\providecommand{\e}{{\mathrm{e}}}
\providecommand{\vect}[1]{\overrightarrow{#1}}
\providecommand{\vectt}[1]{\overrightarrow{\mathrm{#1}}}
\providecommand{\vvt}[1]{\overrightarrow{\mathrm{#1}}}
\providecommand{\np}[2][]{{#2~#1}}
\providecommand{\vv}[1]{{\overrightarrow{#1}}}
\renewcommand{\d}{{\mathrm{d}}}
\providecommand{\degres}{{^\circ}}
\providecommand{\degre}{{^\circ}}
\providecommand{\blue}{\color{blue}}
\providecommand{\red}{\color{red}}
\providecommand{\enskip}{\quad}
\providecommand{\llbracket}{\lbrack\!\lbrack}
\providecommand{\rrbracket}{\rbrack\!\rbrack}
\providecommand{\euro}{{\text{€}}}
\providecommand{\no}{{\text{n°}}}
\providecommand{\interval}[2]{{[#1;#2]}}
\providecommand{\intervFO}[2]{{[#1;#2[}}
\providecommand{\dotfill}{\ldots}

\providecommand{\ds}{\displaystyle}
\providecommand{\textding}[1]{#1}
\providecommand{\Oij}{\(\mathrm{O};\vec{\imath},\vec{\jmath}\)}
\providecommand{\Oijk}{\(\mathrm{O};\vec{\imath},\vec{\jmath},\vec{k}\)}
\providecommand{\Ouv}{\(\mathrm{O};\vec{u},\vec{v}\)}
\providecommand{\intervFF}[2]{\ensuremath{[#1;#2]}}
\renewcommand{\euro}{\ifmmode\text{\texteuro}\else\texteuro\fi}
\DeclareUnicodeCharacter{00B0}{\ifmmode^{\circ}\else\textdegree\fi}
\DeclareUnicodeCharacter{20AC}{\ifmmode\text{\texteuro}\else\texteuro\fi}
% Énoncés : Éducation nationale. Transcription et corrigés : APMEP, Denis Vergès et contributeurs. Découpage et classement : CoopMaths, Sébastien Lozano, Éric Elter et contributeurs. Adaptation de la présentation : Maths Explorer.
% CC BY-SA : https://coopmaths.fr/www/about/
% https://coopmaths.fr/alea/static/bac/2022/tex/bac_2022_05_sujet1_etrangers_1_cor.tex
\begin{document}
\noindent\textit{APMEP / CoopMaths. CC BY-SA. Présentation adaptée par Maths Explorer.}\par\medskip
\medskip

\emph{Cet exercice est un questionnaire à choix multiples. Pour chacune des questions suivantes, une seule des quatre réponses proposées est exacte. Les six questions sont indépendantes.}

\smallskip

\emph{Une réponse incorrecte, une réponse multiple ou l'absence de réponse à une question ne rapporte ni n'enlève de point. Pour répondre, indiquer sur la copie le numéro de la question et la lettre de la réponse choisie.\\
Aucune justification n'est demandée.}

\medskip

\begin{enumerate}
\item On considère la fonction $f$ définie pour tout réel $x$ par $f(x) = \ln \left(1 + x^2\right)$.

Sur $\R$, l'équation $f(x) = \np{2022}$

\begin{center}
\begin{tabularx}{\linewidth}{*{2}{X}}
\textbf{a.~~}n'admet aucune solution. &\textbf{b.~~}admet exactement une solution.\\
\textbf{c.~~}admet exactement deux solutions.&\textbf{d.~~}admet une infinité de solutions.
\end{tabularx}
\end{center}
On a : $f(x) = \np{2022} \iff \ln \left(1 + x^2\right) = \np{2022} \iff \text{e}^{\ln \left(1 + x^2\right)} = \text{e}^{\np{2022}} \iff 1 + x^2 = \text{e}^{\np{2022}} \iff$

$ x^2 = \text{e}^{\np{2022}} - 1$ : ce nombre étant positif, l'équation a deux solutions : $\sqrt{\text{e}^{\np{2022}} - 1}$ et $- \sqrt{\text{e}^{\np{2022}} - 1}$  Réponse \textbf{c.}

\item Soit la fonction $g$ définie pour tout réel $x$ strictement positif par:

\[g(x) = x \ln (x) - x^2\]

On note $\mathcal{C}_g$ sa courbe représentative dans un repère du plan.

\begin{center}
\begin{tabularx}{\linewidth}{*{2}{X}}
\textbf{a.~~} La fonction $g$ est convexe sur $]0~;~+\infty[$.&\textbf{b.~~} La fonction $g$ est concave sur $]0~;~+\infty[$.\\
\textbf{c.~~}La courbe $\mathcal{C}_g$ admet exactement un point d'inflexion sur $]0~;~+\infty[$.&\textbf{d.~~}La courbe $\mathcal{C}_g$ admet exactement
deux points d'inflexion sur $]0~;~+\infty[$.
\end{tabularx}
\end{center}

Sur $\R_{+}$, $g$ est dérivable et sur cet intervalle, $f'(x) = \ln (x) + x \times \dfrac{1}{x} - 2x = \ln (x) - 2x + 1$.

Puis $f''(x) = \dfrac{1}{x} - 2$.

On a donc $f''(x) = 0 \iff \dfrac{1}{x} - 2 = 0 \iff \dfrac{1}{x} = 2 \iff x = \dfrac{1}{2}$.

Sur $]0~;~+ \infty[$, \: $f$ a un seul point d'inflexion d'abscisse $\dfrac{1}{2}$. Réponse \textbf{c.}

\item On considère la fonction $f$ définie sur $]- 1~;~1[$ par

\[f(x) = \dfrac{x}{1 - x^2}\]

Une primitive de la fonction $f$ est la fonction $g$ définie sur l'intervalle $] - 1~;~1[$ par :

\begin{center}
\begin{tabularx}{\linewidth}{*{2}{X}}
\textbf{a.~~} $g(x) = - \dfrac12 \ln \left(1 - x^2\right)$&\textbf{b.~~} $g(x) =  \dfrac{1 + x^2}{ \left(1 - x^2\right)^2}$\\
\textbf{c.~~} $g(x)= \dfrac{x^2}{2\left(x - \dfrac{x^3}{3}\right)}$&\textbf{d.~~} $g(x) = \dfrac{x^2}{2}\ln \left(1 - x^2\right)$
\end{tabularx}
\end{center}

Soit la fonction $g$ définie sur $] - 1~;~1[$ par $g(x) = - \dfrac12 \ln \left(1 - x^2\right)$.

En posant $u(x) = 1 - x^2$, dérivable et non nulle sur $] - 1~;~1[$, on a $g'(x) = - 2x$ et on sait que $g(x) = - \dfrac12\ln u(x)$ entraîne $g'(x) = - \dfrac12\dfrac{u'(x)}{u(x)} = - \dfrac12 \times \dfrac{- 2x}{1 - x^2} = \dfrac{x}{1 - x^2} = f(x)$. Réponse \textbf{a.}
 
\item La  fonction $x \longmapsto  \ln \left(-x^2- x + 6\right)$ est définie sur

\begin{center}
\begin{tabularx}{\linewidth}{*{2}{X}}
\textbf{a.~~} $]- 3~;~2[$&\textbf{b.~~} $]- \infty~;~6]$\\
\textbf{c.~~} $]0~;~+\infty[$&\textbf{d.~~} $]2~;~+\infty[$
\end{tabularx}
\end{center}

La fonction est définie si $- x^2 - x + 6 > 0 \iff x^2 + x - 6 < 0$.

Le trinôme $x^2 +x - 6$ a une racine évidente : 2. Le produit des racines étant égal à $- 6$, l'autre racine est donc $- 3$.

Donc $x^2 + x - 6 = (x - 2)(x + 3)$. On sait que ce trinôme est positif sauf (ce que l'on cherche) entre les racines $- 3$ et 2.

La fonction est donc définie sur l'intervalle $]-3~;~2[$. Réponse \textbf{a.}

\item On considère la fonction $f$ définie sur $]0,5~;~+ \infty [$ par 

\[f(x) =x^2- 4x+ 3 \ln (2x - 1)\]

Une équation de la tangente à la courbe représentative de $f$ au point d'abscisse 1 est:

\begin{center}
\begin{tabularx}{\linewidth}{*{2}{X}}
\textbf{a.~~}$y = 4x- 7$ &\textbf{b.~~} $y = 2x - 4$\\
\textbf{c.~~} $y = -3(x - 1) + 4$ &\textbf{d.~~} $y = 2x - 1$
\end{tabularx}
\end{center}
Une équation de la tangente au point d'abscisse 1 est : $y - f(1) = f'(1)(x - 1)$.

$\bullet~~$$f(1) = 1 - 4 + 3\ln (2 - 1) = - 3 + 3 \times 0 = - 3$ ;

$\bullet~~$$f'(x) = 2x - 4 + 3 \times \dfrac{2}{2x - 1} = 2x - 4 + \dfrac{6}{2x - 1}$, d'où $f'(1) = 2 - 4 + \dfrac{6}{1} = - 2 + 6 = 4$.

Une équation de la tangente est donc $y - (- 3) = 4(x - 1) \iff y = 4x - 4 - 3\iff y = 4x - 7$ . Réponse \textbf{a.}

\item L'ensemble $S$ des solutions dans $\R$ de l'inéquation $\ln (x + 3) < 2\ln (x + 1)$ est:

\begin{center}
\begin{tabularx}{\linewidth}{*{2}{X}}
\textbf{a.~~}$S = ]- \infty~;~-2[ \cup ]1~;~+\infty[$&\textbf{b.~~}$S = ]1~;~+ \infty[$\\
\textbf{c.~~} $S = \emptyset$&\textbf{d.~~} $S = ]- 1~;~1[$
\end{tabularx}
\end{center}
D'après l'énoncé il faut que $x > -3$ et que $x > - 1$. Il faut donc résoudre l'inéquation dans l'intervalle $]- 1~;~+ \infty[$.

$\ln (x + 3) < 2\ln (x + 1) \iff \ln (x + 3) < \ln (x + 1)^2 \iff x + 3 <  (x + 1)^2  \iff  0 < x^2 + 2x + 1 - x - 3 \iff 0< x^2 + x - 2$.

Le trinôme $x^2 + x - 2$ a une racine évidente 1 ; comme le produit des racines est égal à $- 2$, l'autre racine est $- 2$. On a donc :

$x^2 + x - 2  > 0 \iff (x - 1)(x + 2)> 0$ : le trinôme est positif (ce que l'on cherche) sauf entre les racines. D'après la remarque préliminaire  $S = \left]1~;~+ \infty\right[$. Réponse \textbf{b.}
\end{enumerate}

\bigskip

\end{document}
